\subsection{Procedimiento}

Se ejecutaron las cuatro estrategias que pide la guia (UCS, A*+h=0, A*+heuristica basica,
A*+heuristica propuesta) sobre \texttt{tinyCorners}, el layout de referencia adoptado en la
Actividad 7 (ver el hallazgo sobre \texttt{mediumCorners} en esa seccion). La funcion
\texttt{demo\_actividad9()} en \texttt{pacman/searchAgents.py} corre las cuatro por separado (cada una
con su propia instancia de \texttt{CornersProblem}, para que el contador \texttt{\_expanded} de
cada una no se mezcle con las demas) y verifica que las cuatro llegan al mismo costo optimo antes
de reportar los resultados.

\subsection{Resultados}

\begin{table}[H]
\centering
\caption{Experimento comparativo sobre CornersProblem (tinyCorners)}
\label{tab:act9-comparativo}
\begin{tabular}{lrrrr}
\toprule
\textbf{Método} & \textbf{Costo} & \textbf{Expandidos} & \textbf{Tiempo (s)} & \textbf{Óptimo} \\
\midrule
UCS                          & 22 & 295 & 0.002134 & sí \\
A* + h=0                     & 22 & 295 & 0.001887 & sí \\
A* + heurística básica       & 22 & 147 & 0.001370 & sí \\
A* + heurística propuesta    & 22 & 119 & 0.001297 & sí \\
\bottomrule
\end{tabular}
\end{table}

Como es de esperarse, UCS y A*+h=0 expanden exactamente el mismo numero de nodos (295): sin
heuristica, A* se reduce a UCS. Las dos heuristicas disenadas en la Actividad 8 reducen la
exploracion sin perder optimalidad: las cuatro estrategias encuentran el mismo costo (22).

\subsection{Factor de reducción de expansiones}

$$R = \frac{N_{UCS}}{N_{A^*}} = \frac{295}{119} = 2.48$$

usando la heuristica propuesta (la mas informada de las dos disenadas) como referencia de
$N_{A^*}$. Esto significa que UCS expandio aproximadamente 2.48 veces mas estados que
A*+heuristica propuesta para resolver el mismo problema con el mismo costo optimo. Si se usa la
heuristica basica en su lugar, $R = 295/147 = 2.01$: tambien reduce la exploracion a la mitad,
pero menos que la propuesta.

\subsection{Análisis}

La relacion entre la calidad de la heuristica y el numero de nodos expandidos es directa: entre
mas informacion aporte la heuristica sobre el costo real restante (sin dejar de ser una cota
inferior, es decir, sin perder admisibilidad), menos nodos necesita expandir A* para encontrar la
misma solucion optima. Esto se observa con claridad en la progresion
$h=0 \to$ basica $\to$ propuesta: cada heuristica anadida es al menos tan informada como la
anterior (la propuesta siempre devuelve un valor mayor o igual que la basica, porque incluye el
mismo termino de distancia a la esquina mas lejana mas el termino adicional de distancia entre
pares de esquinas pendientes), y en cada paso el numero de nodos expandidos baja: 295 $\to$ 147
$\to$ 119. Ninguna de las dos heuristicas hace que A* pierda optimalidad (las cuatro estrategias
llegan al mismo costo 22) --exactamente lo que garantiza que una heuristica sea admisible--, pero
si reducen sustancialmente cuanto tiene que explorar el algoritmo para encontrarla.
